\scanpage{36}
% archive foliation: 35
la formation de $\underline{\operatorname{Hom}}_{X_i}(F_i,G_i)$ commute à tte extension de base $T'\unsure{-|\to}T_i$.

[pour ceci, écrire une \deleted{résolution} présentation libre de type fini
\[
L_1\longrightarrow L_0-\overline F\longrightarrow0,
\]
\deleted{puis décomposer en
\(0-Z_1-L_1-Z_0-0,\)
et \(0\longrightarrow Z_0\longrightarrow L_0-\overline F-0\)
et prendre les $Z_i$ \unsure{plats puis}}
de sorte qu'on est ramené à regarder
\[
\operatorname{Ker}\bigl(\underline{\operatorname{Hom}}(L_0,F)\longrightarrow\underline{\operatorname{Hom}}(L_1,F)\bigr),
\]
et par un lemme général on peut trouver une décomposition \inserted{de $S'''$} en des $T_i$, telle que, se restreignant à $T_i$, la formation du Ker commute au changement de base quelconque. La conclusion comme plus haut.

\margin{Rédiger sous forme de \unsure{lemme}; les $L_i$ ne \deleted{\unsure{peuvent}} \unsure{être} trouvés que \uline{localement} sur $X$. Il y a lieu de donner un lemme général \unsure{utilisable|valable} en IV, §3.}

Le cas des $\operatorname{Tor}_i$ et $\operatorname{Ext}^i$ est analogue, on est ramené au \unsure{lemme} que voici

\uline{Lemme} Soit $Y$ de type fini sur $T$ noeth, $K$ un complexe de faisceaux cohérents, à op. de degré $+1$ pour fixer les idées, $N$ un entier. Alors il y a une décomposition de $T$ en des $T_i$ loc fermés \inserted{disjoints} telle que
